\documentclass[11pt]{article}
\usepackage[letterpaper,margin=1.1in]{geometry}
\usepackage[T1]{fontenc}
\usepackage{lmodern}
\usepackage{microtype}
\usepackage{amsmath}
\usepackage{booktabs}
\usepackage{array}
\usepackage[round,authoryear]{natbib}
\usepackage{hyperref}
\hypersetup{colorlinks=true,linkcolor=black,citecolor=black,urlcolor=black,
  pdftitle={Ground-to-Grid Scale Factor in the GoSurvey Traverse Editor},
  pdfauthor={Chet Jones}}
\usepackage{enumitem}
\setlength{\parindent}{0pt}
\setlength{\parskip}{0.55\baselineskip}
\newcolumntype{L}[1]{>{\raggedright\arraybackslash}p{#1}}

\title{\textbf{Ground-to-Grid Scale Factor in the\\GoSurvey Traverse Editor}}
\author{Chet Jones \\ \small GoSurvey Project}
\date{October 2026}

\begin{document}
\maketitle
\vspace{-1.5em}
\begin{center}
\small REQ-401 \quad\textbar\quad GitHub Issue \#828 \quad\textbar\quad Pull Request \#833
\end{center}
\vspace{1em}

\begin{abstract}
\noindent GoSurvey's Traverse Editor reduces field-measured slope distances to
horizontal distances, but prior to this change it had no facility for
reducing those horizontal (ground) distances to grid distances under a map
projection. This note derives the governing formulas from standard survey
references, works two examples by hand --- one with an assumed scale
factor and one against a published National Geodetic Survey (NGS)
datasheet --- and compares each hand result against the value the software
actually computes. It also documents the automated test suite that keeps
the implementation honest going forward.
\end{abstract}

\section{What problem this fixes}

A field survey measures \emph{ground distances}: the actual taped or
electronic-distance-measurement (EDM) distance along the earth's curved,
elevated surface. A map projection --- for example MGA56/GDA2020, or a
U.S.\ State Plane Coordinate zone --- represents the earth on a flat
\emph{grid}, which necessarily stretches or shrinks distances by a small,
position-dependent factor $k$ \citep{snyder1987}. A traverse closed and
reported in ground distances will not agree with coordinates taken from the
grid unless every distance is first reduced by $k$ \citep{ghilaniwolf2015}.
GoSurvey's Traverse Editor computed only ground distances; this change adds
an option to reduce them to grid before they drive coordinates, closure,
and the least-squares adjustment.

\section{The governing formulas}

GoSurvey already computes the projection's point scale factor
$k_{\mathrm{grid}}$ through its CS-MAP--based projection engine, and an
elevation (sea-level) reduction factor for the Drawing Settings
Transformation tab. The elevation factor is the standard ratio of the
ellipsoid radius to the ellipsoid radius plus the station's elevation above
it,
\begin{equation}
k_{\mathrm{sea}} \;=\; \frac{R}{R + h},
\label{eq:sea}
\end{equation}
as defined for the U.S.\ State Plane Coordinate System by
\citet{stem1990} and reproduced in standard surveying texts
\citep{ghilaniwolf2015}. The present change reuses that existing pair of
factors rather than introducing a second calculation:
\begin{align}
k_{\mathrm{combined}} &= k_{\mathrm{grid}} \cdot k_{\mathrm{sea}},
  \label{eq:combined}\\
k_{\mathrm{grid\text{-}only}} &= k_{\mathrm{grid}}.
  \label{eq:gridonly}
\end{align}
Equation~\ref{eq:combined}, the product of the grid (projection) scale
factor and the elevation factor, is the \emph{combined factor} of
\citet{ghilaniwolf2015} and \citet{stem1990} --- the standard quantity for
reducing a ground distance measured at elevation to its equivalent grid
distance.

For each traverse leg, the slope distance is first reduced to a horizontal
(ground) distance $d_{\mathrm{ground}}$ by the usual trigonometric leveling
relation \citep{ghilaniwolf2015}. This change inserts one multiplication
before that distance is used for anything else:
\begin{align}
d_{\mathrm{grid}} &= d_{\mathrm{ground}} \cdot k, \\
\Delta E &= d_{\mathrm{grid}} \cdot \sin(\mathrm{bearing}), \qquad
\Delta N = d_{\mathrm{grid}} \cdot \cos(\mathrm{bearing}),
\end{align}
where $k$ is $k_{\mathrm{combined}}$ or $k_{\mathrm{grid\text{-}only}}$ by
the user's choice, and $k = 1$ when the option is left off. The grid
distances $d_{\mathrm{grid}}$ are summed for the closure perimeter and
linear closure, and are the observations handed to the weighted
least-squares traverse adjustment, following the standard treatment of
\citet{ghilani2017}.

\emph{Implementation.} \texttt{src/traverse/TraverseCalc.cpp}, function
\texttt{ComputeTraverse}; the multiplier is resolved by
\texttt{src/commands/CadCommands\_Geo.cpp}, function
\texttt{ResolveTraverseScaleFactor}.

\section{Worked example A: a single leg under a known factor}
\label{sec:exa}

A representative central-meridian grid factor, $k = 0.9996$, applied to a
due-east (bearing $90^{\circ}$) leg measured at $1000.0000$~ft on the
ground:

\begin{center}
\begin{tabular}{@{}llr@{}}
\toprule
Quantity & Formula & Value \\
\midrule
Ground distance $d_{\mathrm{ground}}$ & measured & $1000.0000$ ft \\
Scale factor $k$ & chosen (grid-only or combined) & $0.9996$ \\
Grid distance $d_{\mathrm{grid}}$ & $d_{\mathrm{ground}} \cdot k$ & $999.6000$ ft \\
$\Delta E$ & $d_{\mathrm{grid}} \sin 90^{\circ}$ & $999.6000$ ft \\
$\Delta N$ & $d_{\mathrm{grid}} \cos 90^{\circ}$ & $0.0000$ ft \\
Perimeter contribution & $d_{\mathrm{grid}}$ & $999.6000$ ft \\
\bottomrule
\end{tabular}
\end{center}

This hand result is verified by the test suite: \texttt{tests/TraverseTests.cpp},
\emph{``Scale factor on scales distance, coordinates, and perimeter''}
asserts \texttt{computedHorizDist}, \texttt{computedEasting}, and
\texttt{closurePerimeter} all equal $999.6$ for this input, while
\texttt{computedGroundHorizDist} stays $1000.0$, the unscaled value
preserved for display. A sibling test with the option left off (the
default) shows the output is identical to what the software produced
before this change.

\section{Worked example B: a real projection, checked against an NGS datasheet}
\label{sec:exb}

This example checks GoSurvey's own $k_{\mathrm{grid}}$ at a real,
independently published station rather than an assumed value. The station
is NGS datasheet PID AG9976, ``UNIV OF TEXAS TOWER SEC,'' Travis County,
Texas, on the Texas State Plane Coordinate System, Central Zone
\citep{ngsAG9976}. The datasheet publishes its own computed scale factor
for this station:
\begin{quote}
NGS AG9976, SPC TX-C: $E = 10{,}077{,}391.26$ ftUS, $N = 3{,}115{,}243.14$
ftUS, published scale factor $= 0.99995905$.
\end{quote}
GoSurvey resolves this coordinate through its projection engine
(\texttt{geo::GridScaleFactor}) and through
\texttt{ResolveTraverseScaleFactor}, the function the Traverse Editor
calls, and obtains:

\begin{center}
\begin{tabular}{@{}L{0.46\textwidth}L{0.18\textwidth}L{0.22\textwidth}@{}}
\toprule
Source & $k_{\mathrm{grid\text{-}only}}$ & Agreement \\
\midrule
NGS datasheet (published) & $0.99995905$ & reference \\
\texttt{geo::GridScaleFactor} (direct) & matches to $\le 10^{-8}$ & within $10^{-8}$ \\
\texttt{ResolveTraverseScaleFactor}, Grid only & $0.9999590531$ & within $10^{-7}$ \\
\texttt{ResolveTraverseScaleFactor}, Combined, sea level off & $0.9999590531$ & identical to grid-only \\
\bottomrule
\end{tabular}
\end{center}

\subsection*{Elevation (sea-level) reduction, combined factor}

Turning on the Transformation tab's sea-level option with an elevation
$h = 1000$ ft above an ellipsoid radius $R = 20{,}906{,}000$ ft
($R \approx 6{,}372{,}161.54$ m, converting by the U.S.\ survey foot),
Equation~\ref{eq:sea} gives
\begin{equation}
k_{\mathrm{sea}} = \frac{6{,}372{,}161.544\text{ m}}
{6{,}372{,}161.544\text{ m} + 304.8006\text{ m}} = 0.9999521691.
\end{equation}
GoSurvey computes, through \texttt{geo::SeaLevelScaleFactor} and then the
combined path of \texttt{ResolveTraverseScaleFactor}:

\begin{center}
\begin{tabular}{@{}lrr@{}}
\toprule
Quantity & Hand formula & GoSurvey computed \\
\midrule
$k_{\mathrm{sea}}$ & $0.9999521691$ & $0.9999521691$ \\
$k_{\mathrm{grid}}$ (unchanged by elevation) & --- & $0.9999590531$ \\
$k_{\mathrm{combined}} = k_{\mathrm{grid}} \cdot k_{\mathrm{sea}}$ & $0.9999112247$ & $0.9999112242$ \\
\bottomrule
\end{tabular}
\end{center}

The hand-computed $k_{\mathrm{sea}}$ matches GoSurvey's figure to nine
decimal places, and the hand-computed $k_{\mathrm{combined}}$ matches to
six decimal places (the small residual is rounding in the hand conversion
of $R$ to meters above). $k_{\mathrm{combined}}$ differs from
$k_{\mathrm{grid\text{-}only}}$ exactly because elevation reduction is now
included, confirming that ``Combined'' and ``Grid only'' genuinely produce
two different, individually correct numbers rather than one option doing
nothing.

This is verified by \texttt{tests/TransformationTests.cpp}, test case
\texttt{"ResolveTraverseScaleFactor resolves Combined/GridOnly and reports
why it can't"}, which reuses the same AG9976 fixture as the pre-existing
REQ-360 datasheet test.

\section{The refusal path}

A drawing with no coordinate system assigned --- the default, ``No Datum,
No Projection'' --- must not silently apply an arbitrary number.
\texttt{ResolveTraverseScaleFactor} returns a non-empty reason string and
leaves the factor at exactly $1.0$ for a drawing with an empty
\texttt{zoneCode}. The Traverse Editor shows this reason next to the
checkbox and computes as if the option were off until the drawing's
coordinate system is fixed.

\section{The default (off) path is unchanged}

The same $1000.0000$~ft, bearing $90^{\circ}$ leg from
Section~\ref{sec:exa}, with the option left off (its default state),
produces \texttt{computedHorizDist}, \texttt{computedGroundHorizDist},
\texttt{computedEasting}, and \texttt{closurePerimeter} all equal to
$1000.0000$~ft, exactly what \texttt{ComputeTraverse} produced before this
change, since $d_{\mathrm{grid}} = d_{\mathrm{ground}} \cdot 1.0 =
d_{\mathrm{ground}}$ whenever the option is off.

\section{Automated test summary}

\begin{center}
\small
\begin{tabular}{@{}L{0.27\textwidth}L{0.37\textwidth}L{0.28\textwidth}@{}}
\toprule
File & Test case & What it proves \\
\midrule
\texttt{TraverseTests.cpp} &
  Scale factor off reproduces unscaled coordinates &
  $k=1$ default path is a no-op \\[0.3em]
\texttt{TraverseTests.cpp} &
  Scale factor on scales distance, coordinates, and perimeter &
  $d_{\mathrm{grid}}=d_{\mathrm{ground}}k$ propagates to $\Delta E$, easting, perimeter \\[0.3em]
\texttt{TraverseTests.cpp} &
  Scale factor kind defaults to Combined &
  safe default \\[0.3em]
\texttt{TransformationTests.cpp} &
  ResolveTraverseScaleFactor resolves Combined/GridOnly, reports why it can't &
  no-projection refusal; grid-only matches NGS to $10^{-8}$; combined equals grid-only with sea level off, differs with it on \\
\bottomrule
\end{tabular}
\end{center}

All four cases pass under both the \texttt{GoSurveyTests} and
\texttt{GoSurveySnapTests} executables (Catch2, MSVC release build). The
source is PR~\#833, branch \texttt{feat/req401-traverse-scale-factor},
requirement \texttt{REQ-401} in \texttt{spec/requirements.md}.

\begin{thebibliography}{9}

\bibitem[Ghilani, 2017]{ghilani2017}
Ghilani, C.~D. (2017). \emph{Adjustment Computations: Spatial Data
Analysis}, 6th ed. Hoboken, NJ: John Wiley \& Sons.

\bibitem[Ghilani and Wolf, 2015]{ghilaniwolf2015}
Ghilani, C.~D., and Wolf, P.~R. (2015). \emph{Elementary Surveying: An
Introduction to Geomatics}, 14th ed. Boston: Pearson.

\bibitem[NGS, 2026]{ngsAG9976}
National Geodetic Survey (retrieved 2026). \emph{NGS Data Sheet}, PID
AG9976, ``UNIV OF TEXAS TOWER SEC.'' Silver Spring, MD: NOAA National
Geodetic Survey. \url{https://www.ngs.noaa.gov/}

\bibitem[Snyder, 1987]{snyder1987}
Snyder, J.~P. (1987). \emph{Map Projections --- A Working Manual}. U.S.\
Geological Survey Professional Paper 1395. Washington, DC: U.S.\
Government Printing Office.

\bibitem[Stem, 1990]{stem1990}
Stem, J.~E. (1990). \emph{State Plane Coordinate System of 1983}. NOAA
Manual NOS NGS~5. Rockville, MD: National Geodetic Survey.

\end{thebibliography}

\end{document}
